\documentclass[11pt]{article}
\usepackage[a4paper,margin=25mm]{geometry}
\usepackage[T1]{fontenc}
\usepackage{lmodern}
\usepackage{booktabs}
\usepackage{amsmath}
\usepackage{graphicx}
\usepackage{hyperref}
\usepackage{xcolor}
\usepackage{listings}
\lstset{basicstyle=\ttfamily\small, columns=fullflexible, keepspaces=true,
  breaklines=true, frame=single, framerule=0.3pt, rulecolor=\color{gray}}
\newcommand{\code}[1]{\texttt{#1}}
\newcommand{\us}{\textmu s}

\title{\code{offline.py}: the analysis of one run from the output of the pass}
\author{BGOTSCAN}
\date{2026-10-14}

\begin{document}
\maketitle

\section{Purpose}

The pass of \code{JScan} (\code{jscan\_doc.tex}) reads a run once and
writes, per event, the corrected summed waveform, the charge and time of
every channel in a window, and the single-photoelectron candidates.
\code{offline.py} does the analysis of one run from that file: it
classifies the events, computes the charge of each event in seven ways,
fits the 511~keV photopeak of each, fits the decay of the average pulse,
and writes a per-event table and a summary for the comparison across runs
(\code{scan.py}). It reads only the pass output and the equalization table,
runs in seconds, and is a batch script: no steps, one table on the
terminal, four figures. The definitions of the classes and of the charges
are those of \code{selection\_doc.tex} (steps 6 to 12); this document
gives what is new: the charge noise from the data, the equalized charges,
the photopeak and decay fits, the output.

\section{Usage}

\begin{lstlisting}
python3 src/py/offline.py RUN|FILE [--window A,B] [--smooth 64] [--nsigma 5] [--k 3]
        [--gain FILE] [--dead LIST] [--decay-range A,B] [--save NAME] [--out FILE]
\end{lstlisting}

\begin{center}
\begin{tabular}{l l p{72mm}}
\toprule
option & default & meaning \\
\midrule
\code{RUN|FILE} & & a run name, looked up as \code{pass/RUN\_pass.h5}, or a path \\
\code{--window A,B} & the window of the pass & charge window of the sum: ns before and after the maximum. Set per temperature from the average pulse; the per-channel charges keep the window of the pass \\
\code{--smooth N} & 64 & moving-average width of the pulse search (samples) \\
\code{--nsigma K} & 5 & significance of the charge cut, the pile-up search and the multiplicity \\
\code{--k k} & 3 & zero-suppression threshold of $Q_{\rm kept}$ \\
\code{--gain FILE} & \code{src/spe\_gain\_195K\_41V\_th20.csv} & equalization table of \code{spefit.py}; channels without a factor get 1 and are listed \\
\code{--dead LIST} & the twelve of \code{chanstat\_doc.tex} & channels excluded from every per-channel charge \\
\code{--decay-range A,B} & \code{0.5,4} & \us{} after the maximum over which the exponential is fitted to the average pulse; set per temperature \\
\code{--save NAME} & off & figures saved as \code{src/py/plots/NAME\_classes.png}, \code{NAME\_spectra.png}, \code{NAME\_pulse.png}, \code{NAME\_pileup.png}; otherwise shown \\
\code{--out FILE} & \code{pass/RUN\_offline.h5} & per-event table; the summary goes next to it as \code{RUN\_offline.json} \\
\bottomrule
\end{tabular}
\end{center}

\section{Per event}

$S_i$ is the stored summed waveform divided by its scale; $\Delta t$ = 8~ns;
$a$, $b$ the window in samples from \code{--window}.

\paragraph{Pulse and class.} $\hat S$ = moving average of $S$ over $N$
samples; $i_{\max}$, $H$, $i_{\rm rise}$ as in \code{selection\_doc.tex}
step~7. Charge window $[i_{\max} - a, i_{\max} + b]$ and control window
of the same length in the pulse-free part (step~8); $Q = \Delta t \sum_{\rm
window} S_i$, $Q_{\rm ctrl}$ the same over the control window. The charge
noise is measured on the run: $\sigma_Q$ = rms of $Q_{\rm ctrl}$ over all
events (clipped at three times a first rms, so that pulses in the control
window do not inflate it); this replaces the white-noise formula of the
prototype. $\sigma_N$ = rms of $\hat S$ over the control window, per event.
Class: EMPTY if $Q < K\sigma_Q$; REAL if $Q \ge K\sigma_Q$ and a rising
edge; TAIL if $Q \ge K\sigma_Q$ and no rising edge. Pile-up: the maximum
$P_2$ of $\hat S$ outside the charge window, the valley $V$ between it and
$i_{\max}$; flagged if $P_2 > K\sigma_N$ and $P_2 - V > K\sigma_N$
(step~10); $\Delta t_{\rm pile}$ = time of $P_2$ minus $t_{\max}$.

\paragraph{Charges from the channels.} $q_c$ and $q_c^{\rm ctrl}$ are the
stored per-channel charges in the window of the pass. The charge noise of
channel $c$ is $\sigma_{q,c}$ = clipped rms of $q_c^{\rm ctrl}$ over the
events. Dead channels are removed. The equalization factor $g_c$ comes
from the table; a live channel without one gets $g_c = 1$ and is listed in
the output. With $L$ the live channels,
\[
  Q_{\rm all} = \sum_{c \in L} q_c, \qquad
  Q_0 = \sum_{c \in L,\, q_c > 0} q_c, \qquad
  Q_{\rm kept} = \sum_{c \in L,\, q_c > k\,\sigma_{q,c}} q_c,
\]
and the same three with $q_c / g_c$ in place of $q_c$ (the thresholds on
the unequalized $q_c$, so that the two sets select the same channels):
$Q^{\rm eq}_{\rm all}$, $Q^{\rm eq}_0$, $Q^{\rm eq}_{\rm kept}$. The
multiplicity $M$ is the number of live channels with $q_c > K\sigma_{q,c}$.

\section{Per run}

\paragraph{Classes.} The fraction of EMPTY, REAL and TAIL events, the
fraction of REAL events with pile-up, and the histogram of $\Delta t_{\rm
pile}$ (figure \code{pileup}): a pile-up of random second pulses is flat
in $\Delta t$; a pile of entries just outside the window would be the tail
of the main pulse passing as a second one.

\paragraph{Spectra and photopeak.} For the REAL events without pile-up,
the histogram of each of the seven charges ($Q$ and the six from the
channels), 200 bins from zero to the 99.5th percentile. The photopeak is
located as the rightmost local maximum of the smoothed histogram that
stands above 25\% of its global maximum: the 511~keV peak dominates the
spectrum, the Compton continuum lies below it, and the 1275~keV line is
too weak to pass. It is fitted with a Gaussian on a background,
\[
  f(x) = n\, e^{-(x - \mu)^2 / 2 s^2} + c_0\, \tfrac{1}{2}\operatorname{erfc}\!\big(\tfrac{x - \mu}{s\sqrt 2}\big) + c_1,
\]
the erfc term being the step of the Compton continuum under the peak,
over $\mu \pm 2.5 s$, twice (the range updated after a first fit), by
least squares with Poisson weights. Light yield of the estimator: $\mu$
in ADC\,ns, and in photoelectrons $\mu / q_{\rm pe}$ with $q_{\rm pe} = A_1
\sum_j T_j \Delta t$ the charge of one photoelectron, $A_1$ the median
single-photoelectron amplitude of the gain table and $T$ the template of
the pass; $A_1$ is measured at 41~V only, and the conversion is labelled
with it. Resolution: $\mathrm{FWHM}/\mu = 2.355\, s / \mu$. Figure
\code{spectra}: the seven spectra with their fits; figure \code{classes}:
the $Q$ spectrum of each class.

\paragraph{Average pulse and decay time.} The mean of $S$ over the REAL
events without pile-up, each shifted so that its $i_{\max}$ falls on a
common sample, over the samples all of them cover. A single exponential
$A e^{-t/\tau} + C$ is fitted between $t_{\max} + A$ and $t_{\max} + B$
(\code{--decay-range}); $\tau$ is the decay time of the run. Figure
\code{pulse}: the average pulse and the fit, linear and logarithmic.

\section{Output}

\paragraph{Table.} One line per estimator: position (ADC\,ns and
photoelectrons), error, width, FWHM/position, number of events in the
fit; then the class fractions, the pile-up fraction, $\sigma_Q$, the
decay time, and the channels without an equalization factor.

\paragraph{Per-event file} \code{pass/RUN\_offline.h5}: datasets of length
nev: \code{index}, \code{cls} (0 EMPTY, 1 REAL, 2 TAIL), \code{pileup},
\code{dt\_pile}, \code{t\_max}, \code{H}, \code{t\_rise}, \code{Q},
\code{Q\_ctrl}, \code{Q\_all}, \code{Q\_0}, \code{Q\_kept},
\code{Q\_all\_eq}, \code{Q\_0\_eq}, \code{Q\_kept\_eq}, \code{mult}; the
average pulse as \code{pulse} with its time axis; the options and the
input files as attributes.

\paragraph{Summary} \code{pass/RUN\_offline.json}: the run conditions
copied from the pass (T1, T2, T3, bias, trigger rate), the options, the
class fractions, $\sigma_Q$, the decay time with its error, and for each
estimator the photopeak position, error, width and resolution. This is
what \code{scan.py} reads.

\section{Examples}

\begin{lstlisting}
# the warm run with the pass window
python3 src/py/offline.py 195K_41V_th20 --save off195

# a cold run with a long window and a long decay range
python3 src/py/offline.py 128K_37V_th19 --window 1000,12000 --decay-range 1,10 --save off128
\end{lstlisting}

\end{document}
